\section{Assessing quality of CVs}
Anything coming off an SRV will have
\[
    \myE{\psi_i} = 0 \quad \myE{\psi^2_i} = 1 \quad \psi_i=\sum_j a_{ij}\zeta_j
\]
Work out what this means for when you evaluate the quality of $\psi_i$ when evaluated using ACF.  Will have to validate using full unbiased dataset.

Contemplate the meaning of
\[
    \tau_i \Leftarrow \myE{\psi_i(t+\tau)\psi_i(t)}
\]
and take care about the denominator in the usual ACF.

\subsection{Advances in KDEs and related methods}
From ChatGPT
\begin{enumeratec}
    \item Adaptive bandwidth.  Abramson (1982)
    \[
        \sigma_i \propto\frac{1}{\sqrt{p(x_i)}}
    \]
    Regions that are highly sampled get fine structure.  Silverman (1986) gives low bandwidth will increasing effective sample size.  Generally, bandwidth starts large and gets small.
    \item Anisotropy.  Terrell and Scott (1992).  Boter et al (2010).  Chacon and Duong (2018).
    \item kNN-based.  Choose $\sigma_i$ based on nearest neighbour.
    \item GMM.  NOT KDE.  Karplus, 2009

    See https://pubmed.ncbi.nlm.nih.gov/19284746
    \item Gaussian processes for FES.  Mones, 2016.  ``Exploration, Sampling, And Reconstruction of Free Energy Surfaces with Gaussian Process Regression''
    \item Normalizing flows.  Noe, 2019.  Not a priority for my 2D FES: stuggles with well-separated peaks, overkill for 2D, we lack a priori knowledge of $F(s)$ (Noe leverages Boltzmann factor).
\end{enumeratec}

\subsection{TODO}
\begin{itemize}
    \item review Noe 2019 and look into whether it has a place with a low-dimensional (2D) FES compared to the PES he uses.  I think he leverages
    \[
        p(x) \sim \mye^{-\beta V(x)}
    \]
    to learn the transformation to a $z$ variable that is $3N$ (high) dimensional but with simple structure/shape.  We are trying to learn an FES, so we can similarly leverage knowlege of $F(s)$.
\end{itemize}

\section{Search}
\begin{itemizec}
    \item "2D kernel density estimation" astronomy large dataset
    \item "adaptive kernel density estimation" sharp peaks 2D
    \item "bivariate density estimation" cross-validation bandwidth
    \item "least squares cross validation" kernel density estimation 2D
    \item "spatial hotspot" kernel density estimation bandwidth
    \item "home range estimation" kernel density estimation LSCV
    \item "bivariate KDE" multimodal density estimation
    \item "density estimation" "sharp peaks" "2D"
\end{itemizec}
